\section{均值、相关与平稳性}
\label{sec:06-Random-Processes/01-Foundations/03}

完整的有限维分布族在理论上很完备，但在实际中几乎不可能从数据中完整估计。给你一条长度为 1000 的时间序列，你不可能可靠地估计所有二维、三维、十维联合分布——高维联合分布的直方图需要天文数量的数据才能不稀疏。

工程师和时间序列分析者退而求其次：先看二阶统计量。三个函数概括了过程的前两阶结构：

\[
m_X(t) = \E[X_t],
\qquad
R_X(s,t) = \E[X_s X_t],
\qquad
C_X(s,t) = \Cov(X_s, X_t) = R_X(s,t) - m_X(s)m_X(t).
\]

均值函数 \(m_X(t)\) 描述的是``平均轨迹''——如果我把很多条独立路径叠在一起取平均，得到的那条曲线。自相关函数 \(R_X(s,t)\) 描述两个时刻的原始共同波动——乘积的期望。自协方差函数 \(C_X(s,t)\) 是去均值之后的共同波动——两个时刻各自偏离自身期望多少，以及这些偏离倾向于同号还是反号。

\subsection{自协方差不是可以任意指定的函数}

不是任何看起来光滑、衰减的函数都能充当合法的自协方差。自协方差函数身上背着一条硬约束。

对任意有限个时刻 \(t_1,\ldots,t_n\) 和任意实系数 \(a_1,\ldots,a_n\)，

\[
\sum_{i,j} a_i a_j C_X(t_i, t_j)
= \Var\!\left(\sum_{i=1}^{n} a_i X_{t_i}\right)
\ge 0.
\]

等号来自于协方差的双线性性——将和的方差展开成两两协方差之和，这一步是严格的代数恒等式。不等号来自于方差恒非负——任何随机变量的方差都不可能小于零。

这个约束就是半正定性：对任意有限时刻集，矩阵 \([C_X(t_i, t_j)]_{i,j=1}^{n}\) 必须是半正定矩阵。随便写一个``看起来像相关''的函数——比如在某个滞后上给一个很大的正值——可能产生负的特征值，从而不满足这个约束。自协方差函数的合法空间受限于 Bochner 定理（在宽平稳框架下）或类似的谱表示定理。

另外两条直接的性质有助于日常判断：

\[
C_X(s,t) = C_X(t,s),
\]

自协方差关于时间对称——\(s\) 时刻与 \(t\) 时刻的协方差等于 \(t\) 时刻与 \(s\) 时刻的协方差。以及由 Cauchy--Schwarz 不等式：

\[
|C_X(s,t)| \le \sqrt{C_X(s,s)\,C_X(t,t)}.
\]

相关性不可能超过各自标准差的乘积——相关系数必须在 \([-1,1]\) 之间，这件事在过程中对每一对时刻成立。

\subsection{严平稳：所有联合分布对时间平移不变}

最彻底的平稳性是\textbf{严平稳}（strict stationarity）：把整条过程在时间轴上平移任意量 \(h\)，它的全部概率性质——不仅是均值方差，而是所有有限维联合分布——保持不变。

正式定义：对任意 \(n\)、任意时刻 \(t_1,\ldots,t_n\) 和任意平移 \(h\)，

\[
(X_{t_1}, X_{t_2}, \ldots, X_{t_n}) \overset{d}{=} (X_{t_1+h}, X_{t_2+h}, \ldots, X_{t_n+h}).
\]

这个条件很强。它的直接推论已经包含了实践中常用的全部二阶简化：

\begin{itemize}
\tightlist
\item
  取 \(n=1\)：\(X_t \overset{d}{=} X_{t+h}\) 对任意 \(h\) 成立。所以边缘分布与 \(t\) 无关。若期望存在，\(\E[X_t] = \mu\) 是常数。
\item
  取 \(n=2\)，平移 \(h = -s\)：\((X_s, X_t) \overset{d}{=} (X_0, X_{t-s})\)。所以二维联合分布只依赖时差 \(t-s\)。若二阶矩存在，\(C_X(s,t) = C_X(0, t-s)\)——自协方差只是时差的函数。
\end{itemize}

严平稳推出均值为常数、协方差只依赖时差——前提是这些矩存在。不存在的矩不在讨论范围内。

什么会破坏严平稳？均值随时间漂移（逐年气温上升）、方差随时间变化（金融市场波动率聚集和消散）、或季节效应（月度销售有周期性模式）——任何一种统计特征与绝对时间的挂钩都会破坏平稳性。平稳不要求过程没有波动，只要求波动的统计性质不随时间改变。

\subsection{宽平稳：只要求二阶结构不变}

严平稳的条件涉及所有有限维分布，检验它需要过于丰富的数据。实践中更常使用的是\textbf{宽平稳}（weak stationarity，也叫二阶平稳或协方差平稳）。它只约束一阶和二阶矩：

\begin{enumerate}
\tightlist
\item
  均值为常数：
  \[
  m_X(t) = \mu,\qquad \text{对所有 } t;
  \]
\item
  自协方差只依赖时差：
  \[
  C_X(s,t) = \gamma_X(t-s),\qquad \text{对所有 } s,t.
  \]
\end{enumerate}

第二条需要解释一下记号。原来的 \(C_X(s,t)\) 是一个二元函数——给定两个绝对时间 \(s\) 和 \(t\)，返回一个数。宽平稳条件下，这个二元函数退化为只依赖差值 \(t-s\) 的一元函数。为了避免混淆，通常引入一个新符号 \(\gamma_X(\tau)\)，其中 \(\tau = t-s\) 是滞后（lag），

\[
\gamma_X(\tau) = \Cov(X_t, X_{t+\tau}),\qquad \text{与 } t \text{ 无关}.
\]

这样 \(C_X\) 仍然指二元版本，\(\gamma_X\) 指一元滞后版本，泾渭分明。

宽平稳与严平稳的关系：若过程是严平稳的，且二阶矩有限，则自动满足宽平稳（均值常数和协方差只依赖时差都从严平稳推出来）。反过来，宽平稳一般不蕴含严平稳——高阶联合分布完全可能随时间变化，尽管均值和协方差保持稳定。

Gaussian 过程是这条鸿沟上的桥梁。上一篇已经建立了：Gaussian 过程的有限维分布由均值函数和协方差函数完全决定。因此对于 Gaussian 过程，宽平稳自然推出严平稳——二阶结构锁定了整个分布，平移不变的二阶结构意味着平移不变的全部有限维分布。

\subsection{平稳与平稳增量不是同一回事——随机游走的论证}

平稳增量在上一篇中作为构造条件出现，这里需要和过程本身的平稳性做一次明确对照。

设 \(Z_1, Z_2, \ldots\) 独立同分布，定义随机游走

\[
S_n = Z_1 + Z_2 + \cdots + Z_n,\qquad S_0 = 0.
\]

增量 \(S_{n} - S_{n-1} = Z_n\) 是独立同分布的，因此过程有平稳独立增量：对于不重叠区间，增量独立；每个增量（单步）同分布。

但过程本身不平稳。看方差：

\[
\Var(S_n) = \Var(Z_1 + \cdots + Z_n) = n\Var(Z_1).
\]

方差随 \(n\) 线性增长，不可能是常数——而宽平稳要求方差恒定。随机游走越走越散，水平值的分布随时间膨胀。

反过来，如果过程本身是严平稳的，其增量自动平稳：\((X_{t+h} - X_{s+h})\) 和 \((X_t - X_s)\) 分别是两个时刻对的差分，平移后联合分布不变，所以增量分布相同。注意这里的``差异分布相同''只是增量分布与绝对起点无关；增量分布仍然可以依赖间隔长度。

水平值与增量，两个概念在平稳性讨论中不能互换。说``过程平稳''是在说水平值的统计性质不随时间变化；说``增量平稳''是在说变化量的统计性质与起点无关。一个过程可以拥有平稳增量而不平稳（随机游走），也可以平稳而增量不平稳（不太常见但可能）。

\subsection{平稳不等于遍历——随机水平线的教训}

平稳性保证了总体的统计特征不随时间变化。但你要从数据中估计这些特征，就需要从一条足够长的路径上计算时间平均。平稳性能保证时间平均有意义吗？不保证。

第一篇中构造的随机水平线是良药：

\[
X_t = Z,\qquad Z \sim \mathcal N(0,1).
\]

这个过程是严平稳的——对所有时刻都等于同一个随机常数，任意有限维分布在平移下完全相同。总体均值是 \(\E[X_t] = \E[Z] = 0\)。但任意一条实现——即固定 \(Z=z\) 后——时间平均是

\[
\lim_{T\to\infty} \frac1T \int_0^T X_t\,dt = z,
\]

等于该路径的随机常数，通常不等于零。时间平均收敛到一个随机变量，不是收敛到总体均值。

这就是遍历性（ergodicity）缺失的症状。平稳保证总体分布不随时间变化；遍历保证单条足够长的路径上的时间平均等于总体平均。两个条件各自独立，必须分别检查。

在现实中，遍历性缺失不是一个靠哲学辩论就能绕过的概念性问题。你通常只有一条长时间序列——比如某只股票的过去十年日收益率——而不是同一历史时期在无数平行世界中的样本。如果过程不遍历，仅凭这一条实现无法可靠估计总体统计量。ARMA 建模和谱分析的全部工具默认假设了某种形式的遍历性，但在遇到长记忆过程、体制转换或非平稳趋势时，这些假设需要被逐一复核。

\subsection{去趋势与平稳化——谨慎行事}

面对非平稳序列，建模时的常见第一步是试图把它变成平稳的。两个常用手段：

\begin{itemize}
\tightlist
\item
  \textbf{去趋势}：拟合一条趋势曲线（线性、多项式、局部平滑），从原始序列中减去。
\item
  \textbf{差分}：用 \(X_t - X_{t-1}\) 替代 \(X_t\)，消除随机游走型的非平稳。
\end{itemize}

这些操作的风险在于，它们可能只修好了一个方面却破坏了别的。减去线性趋势可能让均值近似稳定，但方差、自相关结构或季节效应仍可能随时间变化。过度差分（对已经平稳的序列再做差分）会人为引入负自相关并抹掉有价值的低频信息。

建模前值得逐项检查：

\begin{itemize}
\tightlist
\item
  均值是否随时间变化？画滑动平均或局部回归线。
\item
  方差是否稳定？看不同时间窗内的波动幅度。
\item
  自相关是否只依赖滞后而不依赖绝对时间？分段计算自相关函数对比。
\item
  是否存在周期性或结构突变？周期性可以通过谱分析或季节分解识别，突变点需要变点检测方法。
\item
  观测机制本身是否发生过变化？采样频率改变、传感器升级、定义修订——这些不在数据内部，却直接改变数据的统计面貌。
\end{itemize}

``假设平稳''是一个可检验、可失败的工作假设，不是时间序列的默认属性。真正稳健的建模流程是：先假设平稳，用数据检验这个假设，如果不通过则寻找非平稳的根源并相应修正模型——而不是跳过检验直接套用平稳工具。
